% Generated draft: requires mathematical and source review.
\documentclass[11pt]{article}
\usepackage[margin=25mm]{geometry}
\usepackage{fontspec}
\setmainfont{Latin Modern Roman}
\usepackage{amsmath,amssymb,mathtools,unicode-math}
\setmathfont{Latin Modern Math}
\usepackage{enumitem,needspace,xurl,hyperref}
\setlength{\parindent}{0pt}
\setlength{\parskip}{5pt}
\setlength{\emergencystretch}{3em}
\newcommand{\N}{\mathbb{N}}
\newcommand{\Z}{\mathbb{Z}}
\newcommand{\Q}{\mathbb{Q}}
\newcommand{\R}{\mathbb{R}}
\newcommand{\C}{\mathbb{C}}
\newcommand{\F}{\mathbb{F}}
\newcommand{\A}{\mathbb{A}}
\newcommand{\PP}{\mathbb{P}}
\newcommand{\E}{\mathbb{E}}
\newcommand{\eps}{\varepsilon}
\newcommand{\dd}{\,\mathrm{d}}
\begin{document}
\section*{UnsolvedMath ID: 11100037}
\textbf{Source problem code:} catalogue entry
\subsection{Definitions and mathematical statement}
For graded-commutative algebras of characteristic \(p\), a conventional Quillen \(F\)-isomorphism (written ``f-isomorphism'' in the source) is a homomorphism with nilpotent kernel and with power-surjectivity: for each homogeneous element \(b\) of the target, some \(b^{p^r}\) lies in its image. Neither a uniform exponent nor an actual ring isomorphism is required. The source does not state its exact convention.

\emph{Axiomatic stable-homotopy problem 5.} Can Palmieri's methods work in other specific non-Noetherian stable homotopy categories? In particular, following Haynes Miller's suggestion, use them to calculate group cohomology of \(\operatorname{GL}(\infty)\) up to f-isomorphism. A precise version must specify the coefficient/base ring and topology of that group, the cohomology coefficients, and an explicit computable comparison ring and f-isomorphism or appropriate comparison diagram.

For example, only after a ring \(R\) is selected can one define a standard stable linear group \(\operatorname{GL}_\infty(R)=\varinjlim_m\operatorname{GL}_m(R)\), using block stabilization. Discrete group cohomology and a topological-group version are different targets. The excerpt gives none of these choices, and \(\operatorname{GL}(\infty)\) is not silently identified with \(\operatorname{GL}_\infty(\F_p)\).

The source reports that John Palmieri found a computable ring f-isomorphic to the Adams spectral sequence's \(E_2\)-term, describing this as determination up to nilpotence; Hovey qualifies the prime by ``I think this is only at \(p=2\).'' In the usual mod-two sphere setup that term is \(\operatorname{Ext}_{\mathcal A}^{s,t}(\F_2,\F_2)\), but the exact spectral-sequence setting and comparison ring are absent. The report concerns \(E_2\) up to f-isomorphism, not its full ring structure, differentials, abutment, or an explicit computation of sphere homotopy groups. The other categories and the sense in which they are non-Noetherian also require specification.
\subsection{Short English statement}
Extend Palmieri's nilpotence/f-isomorphism method to other non-Noetherian stable settings, especially the cohomology of a specified stable general linear group. The group, coefficients, and comparison ring are missing, and the source's prime-two scope is explicitly tentative.
\subsection{Sources}
\begin{itemize}
\item Local UnsolvedMath registry, permanent ID 11100037, source code catalogue entry; proposer Mark Hovey, ``Axiomatic stable homotopy,'' item 5. John Palmieri is credited for the Adams \(E_2\) comparison and Haynes Miller for the \(\operatorname{GL}(\infty)\) suggestion.
\item \emph{Hovey's Algebraic Topology Problem List}, surviving mirror: \url{https://www-users.cse.umn.edu/~tlawson/hovey/}. Recorded access 2026-07-29; extraction html-ordered-list. The mirror has 65 top-level items across nine category pages.
\item The supplied summary refers to Palmieri's ``Quillen stratification for the Steenrod algebra'' and a memoir on Adams \(E_2\), and mentions Friedlander--Mislin and Duy Nguyen in related detection work. No complete bibliography or exact matching theorem is supplied, and its final discussion is truncated.
\item Supplied status partially\_solved and PARTIAL-PROGRESS. The summary describes the framework as established but the specific stable-GL target as not definitively resolved in its audit. Current primary-literature status and redistribution terms remain NEEDS\_REVIEW.
\end{itemize}
\end{document}
